% TOP_PROBLEM: {"id": 8400036, "title": "O32 — Class numbers of imaginary quadratic orders", "category_id": 3, "category": {"id": 3, "name": "graph_theory", "display_name": "Graph Theory", "description": "Problems involving graphs, networks, and their properties.", "slug": "graph-theory", "order_index": 3, "created_at": "2026-07-31T15:26:25.671Z"}, "status": "partially_solved", "problem_number": "AMR-083-0036", "set_id": 13}
% FIRST_POSTED: 2026-10-01
% ENTRY_KIND: statement_only
% UnsolvedMath ID 8400036; source catalogue code AMR-083-0036
% !TeX program = lualatex
% Definition and sources only; this file is not an LLM solution attempt.
\documentclass[11pt,a4paper]{article}
\usepackage[margin=25mm]{geometry}
\usepackage{fontspec}
\setmainfont{lmroman10-regular.otf}[BoldFont=lmroman10-bold.otf,
 ItalicFont=lmroman10-italic.otf,BoldItalicFont=lmroman10-bolditalic.otf]
\usepackage{amsmath,amssymb,mathtools}
\usepackage{unicode-math}
\setmathfont{latinmodern-math.otf}
\AtBeginDocument{\let\setminus\smallsetminus}
\usepackage{xurl,hyperref}
\hypersetup{colorlinks=true,urlcolor=blue}
\setcounter{secnumdepth}{0}
\setlength{\parindent}{0pt}
\setlength{\parskip}{5pt}
\setlength{\emergencystretch}{3em}
\begin{document}
\section*{O32 — Class numbers of imaginary quadratic orders}\label{8400036}
{\small UnsolvedMath ID 8400036 · AMR-083-0036 · Graph Theory · AMR Open Problem Lists}
\subsection{Definitions and mathematical statement}
Let $D=-d<0$ be a quadratic-order discriminant: $d$ is a positive integer and $D\equiv0$ or $1\pmod4$. Consider the primitive positive definite integral binary quadratic forms
\[
Q(X,Y)=aX^2+bXY+cY^2,\qquad b^2-4ac=D,
\]
where $a,b,c\in\mathbb Z$, $a>0$, and $\gcd(a,b,c)=1$. Two forms are properly equivalent when one is obtained from the other by an integral change of variables with determinant one, that is, by an element of $\mathrm{SL}_2(\mathbb Z)$.

The class number $h(D)$ is the number of these proper equivalence classes. This is finite; equivalently it is the cardinality of the group of invertible ideal classes of the imaginary quadratic order of discriminant $D$. Nonfundamental discriminants are allowed, so the order need not be the full ring of integers of its field.

Given $d$ in binary, can $h(-d)$ be computed exactly by a deterministic classical algorithm whose running time is bounded by a fixed polynomial in $\log(2d)$? The input is promised to give a valid negative discriminant. The output is the class number in binary, not a complete enumeration of form classes.

No factorization of $d$, presentation of a class group, or list of reduced forms is supplied. Enumerating classes is a possible finite approach, but a positive answer must establish the requested bit-complexity bound uniformly as the discriminant grows.
\subsection{Short English statement}
Can imaginary quadratic order class numbers be computed exactly and deterministically in time polynomial in the discriminant’s binary length?
\subsection{Sources}
\begin{itemize}
\item[S1] ULAM.AI and the UnsolvedMath Contributors, supplied problems.json, record 8400036 (AMR-083-0036), AMR Open Problem Lists. Catalogue identity and source transcription; CC BY 4.0.
\url{https://huggingface.co/datasets/ulamai/UnsolvedMath}.
\item[S2] Adleman \& McCurley - Open Problems in Number Theory Complexity II (1994), O32, p21 / LNCS p311. Original source indexed by AMR-083-0036.
\url{https://mccurley.org/papers/open.ps.gz}.
\end{itemize}
\end{document}
